\section{Conclusion and Future Work}
\label{sec:conclusion}

In this paper, we have adopted a critical methodological lens to investigate Quantum Wavefunction Superposition Merging (QWS-Merge), a recent state-of-the-art dynamic model merging protocol. By applying Occam's razor, we analyzed the structural and mathematical properties of QWS-Merge and exposed key insights regarding the layer-wise specialization confounder, scale regularization, Softmax zero-sum competitive bottlenecks, and paradigm differences.

To control for these variables and establish rigorous baselines, we proposed the Bounded Classical Router (BC-Router) framework with classical variants: the Bounded Linear Router (BL-Router), the Global Router with Layer-wise Scaling (GLS-Router), and the Softmax-free Bounded Sigmoidal Router (BSigmoid-Router). Our empirical results on a Vision Transformer backbone fine-tuned to true convergence across four distinct image datasets (MNIST, FashionMNIST, CIFAR-10, SVHN) demonstrate that:
\begin{itemize}
    \item While the joint accuracies of QWS-Merge and the unsupervised online test-time adaptation baseline AdaMerging (TTA) are close ($83.97 \pm 0.53$\% vs.\ $89.30 \pm 0.00$\%), they represent completely different operational paradigms. QWS-Merge represents a highly efficient offline calibration approach with zero inference-time latency or active optimization overhead. We provide explicit calibration and inference-time latency benchmarks in Appendix~\ref{sec:appendix_complexity} and Table~\ref{tab:complexity_comparison}, demonstrating that our proposed BSigmoid-Router is over 25 times faster than AdaMerging during inference, thereby resolving the severe latency bottleneck of test-time optimization.
    \item Applying standard L2 regularization to a simple classical routing head completely resolves the SVHN collapse, boosting SVHN accuracy to \textbf{91.73} $\pm$ \textbf{3.71}\% (outperforming QWS-Merge by \textbf{+12.00\%}) and increasing joint homogeneous accuracy from $78.13 \pm 2.54$\% to \textbf{82.80} $\pm$ \textbf{4.85}\%. This proves that classical linear weight-routing is a highly competitive paradigm when properly regularized.
    \item Our proposed GLS-Router controls for layer-wise capacity, but exhibits high sensitivity to calibration sampling across seeds (SVHN std.\ dev.\ of \textbf{24.30\%}). Moreover, unregularized GLS-Router suffers from severe overfitting and collapses on FashionMNIST ($64.80 \pm 3.53$\%). This collapse serves as a major scientific finding, proving that QWS-Merge's wave phase-projection equations function as an essential and highly robust mathematical regularizer that protects dynamic routing from overfitting under extremely tight offline calibration budgets.
    \item Deconstructing the BL-Router collapse resolves a critical logical contradiction, showing that its failure on SVHN ($31.73 \pm 16.03$\%) is not caused by the scale ceiling itself (as static Uniform Merge at $0.3$ scale gets $77.60 \pm 0.00$\%), but is instead a routing head optimization failure on the small calibration set. This is further resolved by our Softmax-free \textbf{BSigmoid-Router}, which eliminates the zero-sum competitive bottleneck to achieve a stable homogeneous joint accuracy of \textbf{83.73} $\pm$ \textbf{1.93}\% and heterogeneous stream accuracy of \textbf{83.96} $\pm$ \textbf{2.27}\% ($B=1$). Remarkably, BSigmoid-Router's extreme parameter efficiency (requiring only 772 parameters) coupled with its feed-forward design makes it a highly practical, plug-and-play solution for real-world Edge AI systems where inference latency and computation budget constraints are paramount.
\end{itemize}

Our work serves as a broader warning and a call to action for the machine learning community. The pressure to publish novel, complex architectures often leads researchers to introduce over-engineered metaphors instead of critically examining and properly optimizing simple classical baselines. We urge the community to return to scientific rigor: prioritize fair and exhaustive baseline tuning, control for key structural confounders, and recognize that identifying the true drivers of performance—such as proper regularization—is just as valuable as proposing new solutions.

\textbf{Limitations \& Future Work:} While our work successfully demystifies dynamic weight-space routing on compact Vision Transformers, we must explicitly acknowledge the scope of our evaluation. Specifically, our empirical sandbox is restricted to a compact backbone (\texttt{vit\_tiny\_patch16\_224}) evaluated on four vision datasets. While this environment is standard, highly controlled, and mathematically ideal for isolating and deconstructing parameter-routing mechanisms without confounding scaling factors, verifying whether these findings generalize to larger model architectures (e.g., Swin Transformers, ViT-Base/Large, or larger multimodal models like CLIP and LLaVA) and larger-scale datasets (such as ImageNet classification sub-tasks, DomainNet, or multi-domain benchmarks) remains a prioritized path for future empirical validation. Furthermore, future work should explore whether these classical scale-regularization, calibration controls, and Softmax-free independent sigmoidal routing insights generalize to billion-parameter large language models (such as LLaMA-1B/3B) across layer-wise attention and MLP blocks, especially under sparse gating topologies (e.g., Top-1 or Top-2) where task suites scale to $K \ge 10$ experts.
